\chapter{2023年「高等数学A1」\titlelogo{https://sci.hdu.edu.cn}{HDU}期中回忆}

\date{2023年12月2日}{基础数学教学团队}{A0714201}
{未央数学讲师~llf}
{\qrcode{https://qm.qq.com/q/OlkSiy9W04}}{\href{https://qm.qq.com/q/OlkSiy9W04}{HDU 数学营}}
{\qrcode{http://weixin.qq.com/r/mhKYgFTEMQlOrRDj90eI}}{\href{http://weixin.qq.com/r/mhKYgFTEMQlOrRDj90eI}{未央学社公众号}}

\section{选择题（每题3分，共36分）}
\begin{choice}{D}
    $\lim_{x\to 1}\arctan\frac{1}{x-1}=$
\begin{tasks}(4)
    \task $\frac\pi2$
    \task $-\frac\pi2$
    \task $\infty$
    \task 不存在
\end{tasks}
\end{choice}
\begin{solution}
\begin{itemize}
    \item $x\to 1^+$时$\frac{1}{x-1}\to +\infty,\ \lim_{x\to 1}\arctan\frac{1}{x-1}=\frac\pi2$.
    \item $x\to 1^-$时$\frac{1}{x-1}\to -\infty,\ \lim_{x\to 1}\arctan\frac{1}{x-1}=-\frac\pi2$. 所以极限不存在. \sokka{D}
\end{itemize}
\end{solution}

\begin{choice}{D}
    设$\alpha_1=1-\cos x,\ \alpha_2=\sqrt x\ln(1+x),\ \alpha_3=\e^{\sqrt x}-1,\ \alpha_4=\sqrt[3]{1+x}-1$，则$x\to 0^+$时以上四个无穷小量按照从低阶到高阶的排序是
\begin{tasks}(4)
    \task $\alpha_4,\ \alpha_3,\ \alpha_2,\ \alpha_1$
    \task $\alpha_3,\ \alpha_2,\ \alpha_4,\ \alpha_1$
    \task $\alpha_3,\ \alpha_4,\ \alpha_1,\ \alpha_2$
    \task $\alpha_3,\ \alpha_4,\ \alpha_2,\ \alpha_1$
\end{tasks}
\end{choice}
\begin{solution}*
    $1-\cos x\sim\frac12x^2,\ \sqrt x\ln(1+x)\sim x^{\frac32},\ \e^{\sqrt x}-1\sim\sqrt x,\ (1+x)^{\frac13}-1\sim\frac13 x$. \sokka{D}
\end{solution}

\begin{choice}{D}
    已知$y=\ln x+\ln\ln x$，则$\d y/\d x=$
\begin{tasks}(4)
    \task $\frac{1}{x\ln x}$
    \task $\frac{1+x}{x\ln x}$
    \task $\frac{x+\ln x}{x\ln x}$
    \task $\frac{1+\ln x}{x\ln x}$
\end{tasks}
\end{choice}
\begin{solution}*
    $(\ln x)'=\frac1x,\ (\ln\ln x)'=\frac{1}{\ln x}\cdot\frac1x=\frac{1}{x\ln x}$. \sokka{D}
\end{solution}

\begin{choice}{B}
    下列各式中正确的是
\begin{tasks}(4)
    \task $\lim_{x\to\infty}\frac{\sin x}{x}=1$
    \task $\lim_{x\to\infty}\frac{\sin x}{x}=0$
    \task $\lim_{x\to0}\frac{\sin x}{x}=0$
    \task $\lim_{x\to\infty}\frac{\tan x}{x}=1$
\end{tasks}
\end{choice}
\begin{solution}*
    根据第一个重要极限和等价无穷小可得.
\end{solution}

\begin{choice}{B}
    下列命题正确的是
    \begin{tasks}(2)
        \task 无界数列必趋于无穷大
        \task 趋于无穷大的数列必无界
        \task 不趋于无穷大的数列必有界
        \task 趋于正无穷大的数列必在n充分大时单调
    \end{tasks}
\end{choice}
\begin{solution}*
    A选项反例：$y=x\sin\frac1x$为振荡函数. \sokka{B}
\end{solution}

\begin{choice}{A}
    偶函数$f(x)$在实数域$\mathbb{R}$上可导，则$f'(x)$是
    \begin{tasks}(4)
        \task 奇函数
        \task 偶函数
        \task 非奇非偶函数
        \task 不能确定
    \end{tasks}
\end{choice}
\begin{solution}*
    求$n$次导，变换$n$次奇偶性. \sokka{A}
\end{solution}
    
\begin{choice}{C}
    已知函数$f(x),\ g(x)$在$[a,b]$上连续可导，$f(x)\neq 0,\ g(x)\neq 0,\ f'(x)g(x)-f(x)g'(x)<0$，当$x\in(a,b)$时，有
\begin{tasks}(4)
    \task $f(x)g(x)<f(a)g(a)$
    \task $f(x)g(x)<f(b)g(b)$
    \task $\frac{f(x)}{g(x)}<\frac{f(a)}{g(a)}$
    \task $\frac{f(x)}{g(x)}<\frac{f(b)}{g(b)}$
\end{tasks}
\end{choice}
\begin{solution}
\begin{itemize}
    \item 由于$f'(x)g(x)-f(x)g'(x)<0$，所以$\ab(\frac{f(x)}{g(x)})'<0$，$\frac{f(x)}{g(x)}$单调递减.
    \item 由于$a<x<b$，所以$\frac{f(x)}{g(x)}<\frac{f(a)}{g(a)}$. \sokka{C}
\end{itemize}
\end{solution}

\begin{choice}{D}
    下列命题中正确的是
    \begin{tasks}(2)
        \task 极小值一定大于极大值
        \task 极值点一定是驻点
        \task 驻点一定是极值点
        \task 可导函数的极值点一定是驻点
    \end{tasks}
\end{choice}
\begin{solution}
    \begin{itemize}
        \item 极值点可以是驻点、端点，所以B选项错误.
        \item 极小值与极大值无可比性，\emph{\textbf{区分最值与极值的区别.}} 所以A选项错误.
        \item 若二阶导数为零或不存在则可能是拐点，所以C选项错误. \sokka{D}
    \end{itemize}
\end{solution}

\begin{choice}{B}
    如果函数$f(x)$在点$x_0$处不连续，则函数$f(x)$在点$x_0$处
\begin{tasks}(4)
    \task 必可导
    \task 必不可导
    \task 不一定可导
    \task 一定可微
\end{tasks}
\end{choice}
\begin{solution}*
    可导必连续，不连续必不可导.
\end{solution}

\begin{choice}{C}
    已知$y=\ab(\sin x^2)^2$，则$\d y/\d x=$
\begin{tasks}(4)
    \task $4\sin(2x^2)$
    \task $4x\cos(2x^2)$
    \task $2x\sin(2x^2)$
    \task $2x\cos(2x^2)$
\end{tasks}
\end{choice}
\begin{solution}*
    复合函数求导$y'=2\sin x^2\cdot\cos x^2\cdot 2x=2x\sin(2x^2)$. \sokka{C}
\end{solution}

\begin{choice}{C}
    已知当$x\to 0$时，$\e^x-(1+ax+bx^2)$是比$x^2$的高阶无穷小，则
\begin{tasks}(4)
    \task $a=1,\ b=1$
    \task $a=1,\ b=-1$
    \task $a=1,\ b=1/2$
    \task $a=1/2,\ b=1$
\end{tasks}
\end{choice}
\begin{solution}*
    意味着$x\to 0$时$\e^x$可用$1+ax+bx^2$逼近，近似误差要小于等于$x^2$的某个常数，根据$x=0$的泰勒展开有$\e^x=1+x+\frac{x^2}{2!}+\frac{x^3}{3!}+\cdots$. 所以$a=1,\ b=\frac12$.
\end{solution}

\begin{choice}{D}
    已知$\lim_{x\to 0}\frac{f(x)-f(0)}{x^2}=1$，则函$f(x)$在$x=0$处
\begin{tasks}(4)
    \task 可导，且$f'(x)\neq 0$
    \task 不可导
    \task 取得极大值
    \task 取得极小值
\end{tasks}
\end{choice}
\begin{solution}*
    由于$[f(x)-f(0)]/x^2=1>0$，所以$x\to 0$时$f(x)>f(0)$，$f(0)$取极小值.
\end{solution}

\section{填空题（每题3分，共18分）}
\begin{problem}
    极限$\lim_{x\to 0}\frac{\sqrt{1+5x}-\sqrt{1+x}}{x^2+4x}=$\ans{$\frac12$}.
\end{problem}
\begin{solution}*
    $\lim_{x\to0}\frac{1+5x-1-x}{(x^2+4x)\ab(\sqrt{1+5x}+\sqrt{1+x})}=\lim_{x\to 0}\frac{4}{(x+4)\ab(\sqrt{1+5x}+\sqrt{1+x})}=\frac12$.
\end{solution}

\begin{problem}
    曲线$y=2\cos x$在点$(\pi/3,1)$处的曲率是\ans{$1/8$}.
\end{problem}
\begin{solution}*
    代入曲率公式得$k=\frac{\abs{y''}}{(1+y')^{\frac32}}=\frac{\abs{-2\cos\pi/3}}{\ab[1+(-2\sin\pi/3)^2]^{\frac32}}$
\end{solution}

\begin{problem}
    极限$\lim_{x\to 0}(1+3x)^{\frac{1}{\sin x}}=$\ans{$\e^3$}.
\end{problem}
\begin{solution}*
    $\lim_{x\to 0}(1+3x)^{\frac{1}{3x}\cdot\frac{3x}{\sin x}}=\e^3$.
\end{solution}

\begin{problem}
    设参数方程
    $\begin{cases}
        x=1+\frac12t^2\\
        y=\e^{3t}
    \end{cases}$
    ，则$\odv yx\Big|_{t=3}=$\ans{$\e^9$}.
\end{problem}
\begin{solution}*
    $\odv xt=t,\ \odv yt=3\e^{3t}$，所以$\odv yx=\frac{3\e^9}{3}=\e^9$.
\end{solution}

\begin{problem}
    已知$y=\ab(\frac{x}{1+x})^x\ (x>0)$，则$\odv yx=$\ans{$\e^{x[\ln x-\ln(1+x)]}\cdot\ab[\ln x-\ln(1+x)+\ab(1-\frac{x}{1+x})]$}.
\end{problem}
\begin{solution}*
    $y=\e^{x[\ln x-\ln(1+x)]}$，所以$\odv yx=\e^{x[\ln x-\ln(1+x)]}\cdot\ab[\ln x-\ln(1+x)+\ab(1-\frac{x}{1+x})]$.
\end{solution}

\begin{problem}
    曲线$y=x\e^x$的拐点为\ans{$(-2,-2\e^{-2})$}.
\end{problem}
\begin{solution}*
    $y'=\e^x+x\e^x,\ y''=\e^x+\e^x+x\e^x=2\e^x+x\e^x=0,\ x=-2$. 所以拐点为$(-2,-2\e^{-2})$.
\end{solution}

\section{计算题（共46分）}
\begin{problem}*[][10]
    设函数$f(x)=\begin{cases}
        x^2, & x\leq 0\\
        x^2\ln x, & x>0
    \end{cases}$，问
\begin{enumerate}
    \item $f(x)$在$x=0$处是否可导.
    \item 求导函数$f'(x)$得表达式.
    \item 讨论导函数$f'(x)$在$x=0$处是否连续.
\end{enumerate}
\end{problem}
\begin{solution}
\columnseprule=.5pt
\begin{multicols*}{2}
\begin{enumerate}
    \item $\begin{cases}
        f_-'(0)=\lim_{x\to 0^-}\frac{f(x)-f(0)}{x-0}=\lim_{x\to 0^-}\frac{x^2-0}{x}=0\\
        f_+'(0)=\lim_{x\to 0^+}\frac{f(x)-f(0)}{x-0}=\lim_{x\to 0^+}\frac{x^2\ln x}{x}=0
    \end{cases}$. 所以$f_-'(0)=f_+'(0)$，$f(x)$在$x=0$处可导.
    \item 由(1)可知$f'(0)=0$.
    
    $f'(x)=\begin{cases}
        2x, & x<0\\
        0, & x=0\\
        2x\ln x+x, &x>0
    \end{cases}$.
    \item $\lim_{x\to 0^-}f'(x)=f'(0)=0$，所以$f'(x)$在$x=0$连续.
\end{enumerate}
\end{multicols*}
\end{solution}

\begin{problem}*[][10]
    求函数
    $f(x)=\begin{cases}
        x\e^{\frac{1}{x-1}}, &x\geq 0\\
        \arctan\frac1x, &x<0
    \end{cases}$
    的间断点，并判断该间断点的类型.
\end{problem}
\begin{solution}
\begin{itemize}
    \item $\lim_{x\to 1^+}f(x)=+\infty,\ \lim_{x\to 1^-}=0$. 所以在$x=1$处为第二类间断点，且为无穷间断点.
    \item $\lim_{x\to 0^+}f(x)=0\neq\lim_{x\to 0^-}=-\frac\pi2$. 所以在$x=0$处为第一类间断点，且为跳跃间断点.
\end{itemize}
\end{solution}

\begin{problem}*[][10]
    设函数$y=f(x)$由方程$y+x\e^y=1$确定，求$y'(0)$和$y''(0)$的值.
\end{problem}
\begin{solution}
    由题意知当$x=0$时，$y(0)=1$. 两边分别对$x$求两次导得\[y'+\e^y+x\e^yy'=0\ (1)\quad y''+\e^yy'+\e^yy'+x\e^y(y')^2+x\e^yy''=0\ (2)\]
\begin{itemize}
    \item 将$x=0,\ y(0)=1$代入(1)得$y'(0)=-\e$.
    \item 将$x=0,\ y(0)=1,\ y'(0)=-\e$代入(2)得$y''(0)=2\e^2$.
\end{itemize}
\end{solution}

\begin{problem}*[][10]
    求函数$y=1+36x/(x+3)^2$的单调区间与极值，并说明是极大值还是极小值.
\end{problem}
\begin{solution}
    易知定义域为$(-\infty,-3)\cup(-3,+\infty)$，且$y'=\frac{36(x+3)^2-36x\cdot 2(x+3)}{(x+3)^4}=\frac{36(3-x)}{(x+3)^3}$.
\begin{itemize}
    \item 由$y'>0$得$-3<x\leq 3$，单调递增区间：$(-3,3]$.
    \item 由$y'<0$得$x\geq 3\text{ or }x<-3$，单调递减区间：$(-\infty,-3)\cup[3,+\infty)$.
    \item 当$x=3$时，$y$取得极大值为$y(3)=4$，无极小值.
\end{itemize}
\end{solution}

\begin{problem}*[][6]
    函数$f(x)$在$[0,2]$内三阶可导,且$f(0)=f(2)=0$,令$F(x)=x^2f(x)$. 证明$F'''(x)$在$(0,2)$内至少有一个零点.
\end{problem}
\begin{solution}
\begin{proof}*
\begin{itemize}
    \item $f(0)=f(2)=0$，所以$F(0)=F(2)=0$. 由罗尔定理知$\exists\xi_1\in(0,2),\ \text{s.t.}\ F'(\xi_1)=0$.
    \item $F'(x)=2xf(x)+x^2f'(x)$，所以$F'(0)=F'(\xi_1)=0$. 由罗尔定理知$\exists\xi_2\in(0,\xi_1),\ \text{s.t.}\ F''(\xi_2)=0$.
    \item $F''(x)=2f(x)+2xf'(x)+2xf'(x)+x^2f''(x)$，所以$F''(0)=F''(\xi_2)=0$.
    \item 由罗尔定理知$\exists\xi_3\in(0,\xi_2),\ \text{s.t.}\ F'''(\xi_3)=0$. 由于$\xi_3\in(0,\xi_2)\subseteq(0,2)$，所以$\xi_3$是$F'''(x)$在$(0,2)$内的一个零点，$F'''(x)$在$(0,2)$内至少有一个零点.\QED
\end{itemize}
\end{proof}
\end{solution}


\vfill
\rightline{/*Powered by \pkg{LiteSolution} \textsf\LaTeX~Template. Designed by WeYoung \href{https://ctan.org/home/xiamyphys}{Axia}.*/}